\documentclass[12pt]{article}
\usepackage{amsmath}
\usepackage{amssymb}
\usepackage{amsthm}
\usepackage{listings}
\usepackage{xcolor}
\usepackage{hyperref}
\hypersetup{
    colorlinks=true,
    linkcolor=black,
    urlcolor=blue,
    citecolor=blue,
    filecolor=black
}
\usepackage{geometry}
\geometry{letterpaper, margin=1in, top=0.5in, bottom=0.5in}


\newtheorem{theorem}{Theorem}[section]
\newtheorem{definition}{Definition}[section]

\title{The Poison Language: Working Notes}
\author{Vadim Zaliva}
\date{\today}

\lstdefinelanguage{poison}{
  language=caml,
  morekeywords={fn,pfn,unit},
  comment=[s]{(*}{*)},
  commentstyle=\color{green!60!black},
  keywordstyle=\color{blue},
  stringstyle=\color{red},
  showstringspaces=false
}

\lstset{
    basicstyle=\ttfamily\small,
    breaklines=true,
    frame=single,
    numbers=left,
    numberstyle=\tiny,
    keywordstyle=\color{blue},
    commentstyle=\color{green!60!black},
    stringstyle=\color{red},
    showstringspaces=false
}

\begin{document}

\maketitle

\begin{abstract}
    This is a working document. For complete formal definitions, see the
    accompanying PDF document generated from the Ott specification.
\end{abstract}

\tableofcontents

\sloppy

\section{Introduction}
\label{sec:introduction}

We start from the textbook two-point information-flow lattice
($\mathrm{clean} < \mathrm{poisoned}$) and repurpose it to model data-fault
propagation. On top of this we add four ingredients that,
together, are not present in prior work:
\begin{enumerate}
    \item Region-parameterised reference types that distinguish pointer
    reliability from payload reliability by giving every closure or
    variable two orthogonal poison flags.
    \item A flow-sensitive typestate map that is updated after every assignment
    and joined at control-flow merge-points (à la \cite{DBLP:conf/popl/HuntS06})
    so both explicit and implicit flows of faults are tracked.
    \item A predicate system that allows selective de-poisoning based on runtime
    conditions.
    \item Support for higher-order functions with explicit type annotations that capture their body and typing
    context at definition time, then re-type-check the body at application time with the actual argument's
    predicates and current logical environment, enabling predicate-based
    de-poisoning within function bodies while ensuring fault taint cannot be
    hidden behind lambda abstraction.
\end{enumerate}

Together, these features give a single, uniform type system that provides a
fail-safe-integrity guarantee: any fault -- whether injected by a corrupted
sensor value, a poisoned pointer, or a control-flow decision -- propagates
conservatively through both data and control paths, even in the presence of
aliasing and higher-order code.



\section{Language Overview}

A simple imperative language with references and functions:

\begin{lstlisting}[language=poison]
e ::= n | b | x | l                    (* values and variables *)
    | e1 op e2                          (* arithmetic/comparison *)
    | unary_op e                        (* unary operations *)
    | if e1 then e2 else e3             (* conditionals *)
    | ref e | e1 := e2 | !e             (* references *)
    | unit | e1 ; e2                    (* sequencing *)
    | while e1 do e2                    (* loops *)
    | let x = e1 in e2                  (* local bindings *)
    | fn x : ST1 => e : ST2 | e1 e2       (* functions with explicit return types *)
    | pfn x : ST1 => pe : ST2 | pe        (* pure functions with explicit return types *)

pe ::= n | b | x | l                   (* pure expressions *)
     | pe1 op pe2 | if pe1 then pe2 else pe3
     | unary_op pe
     | !pe | unit
     | let x = pe1 in pe2 | pe1 pe2
     (* Note: no ref, :=, while, fn, pfn, or ; allowed *)

(* surface types for annotations *)     
ST ::= int | bool | unit | ref ST | ST1 -> ST2  

(* enriched types *)
T ::= int | bool | unit 
    | ref T @ l         (* region types with location indices *)
    | ST1 -> ST2 # p typing_closure tcl (* impure function types with surface parameter and return types and typing closure *)
    | ST1 -> ST2 # p typing_pclosure ptcl (* pure function types with pure typing closure *)

(* qualified types with poison status and predicate *)    
Q ::= T :> p { P }      

P ::= guard k pred_closure (x, pe, r)  (* predicate as pure closure *)
    | guard k                          (* guard tag only *)
    | no_pred                          (* trivial predicate *)
\end{lstlisting}

\subsection{Three-Level Type System}

The Poison language employs a three-level type hierarchy to separate concerns
between user annotations, internal type system operations, and fault tracking:

\subsubsection{Surface Types (\texttt{ST})}
\textbf{Purpose:} User-facing type annotations for function signatures.
\begin{itemize}
    \item \textbf{Ott notation:} \texttt{styp, ST}
    \item \textbf{Lean implementation:} \texttt{Typ}
    \item \textbf{Grammar:} \texttt{ST ::= int | bool | unit | ref ST | ST1 -> ST2}
    \item \textbf{Usage:} Written by programmers in function parameter and return type annotations
    \item \textbf{Examples:} \texttt{fn x : int => x + 1 : int}, \texttt{pfn f : int -> bool => f 42 : bool}
\end{itemize}

Surface types are the \emph{interface} between the programmer and the type
system. They are deliberately simple and free of implementation details like
region annotations, poison flags, or typing closures. Users write surface types
in their code but never see the more complex internal representations.

\subsubsection{Enriched Types (\texttt{T})}
\textbf{Purpose:} Internal type system representation with region information and typing closures.
\begin{itemize}
    \item \textbf{Ott notation:} \texttt{typ, T}
    \item \textbf{Lean implementation:} \texttt{ETyp} (Enriched Type)
    \item \textbf{Grammar:} \texttt{T ::= int | bool | unit | ref T @ l | ST1 -> ST2 \# p typing\_closure tcl}
    \item \textbf{Usage:} Used internally by the type checker for sound reasoning about references and functions
    \item \textbf{Key additions:}
    \begin{itemize}
        \item \textbf{Region annotations} (\texttt{@l}): Track which specific memory location each reference points to
        \item \textbf{Typing closures} (\texttt{typing\_closure tcl}): Capture function bodies and variable contexts for re-typechecking at application time
        \item \textbf{Function poison} (\texttt{\#p}): Track the poison status of the function itself (from the program counter at definition time)
    \end{itemize}
\end{itemize}

Enriched types extend surface types with the machinery needed for sound type
checking. \emph{Once assigned, an enriched type is \emph{immutable} during
analysis—only the environments grow around it.} The region annotation
\texttt{ref T @ l} solves the soundness problem in dereference by statically
tracking which location each reference targets. Function types \texttt{ST1 ->
ST2 \# p typing\_closure tcl} store both the user-provided surface types and the
implementation details needed for higher-order predicate-based de-poisoning.

\subsubsection{Qualified Types (\texttt{Q})}
\textbf{Purpose:} Complete type information including poison status and predicates for fault tracking.
\begin{itemize}
    \item \textbf{Ott notation:} \texttt{qualified\_type, Q}
    \item \textbf{Lean implementation:} \texttt{QTyp}
    \item \textbf{Grammar:} \texttt{Q ::= T :> p \{ P \}}
    \item \textbf{Components:}
    \begin{itemize}
        \item \textbf{Base type} (\texttt{T}): The enriched type describing the value's structure
        \item \textbf{Poison status} (\texttt{p}): Whether the value is \texttt{clean} or \texttt{poisoned}
        \item \textbf{Predicate} (\texttt{P}): Optional guard for conditional de-poisoning
    \end{itemize}
    \item \textbf{Usage:} All values in typing environments (\texttt{$\Gamma$}, \texttt{$\sigma$}) have qualified types
\end{itemize}

Qualified types are the \emph{working currency} of our typestate system. Every
typing judgement produces a qualified type that captures not just what kind of
value we have, but also its reliability (poison status) and any conditions under
which it might be considered reliable (predicates). \underline{The $(p,P)$
decoration is \emph{mutable}: typing rules may rewrite it as faults propagate,
while the embedded $T$ remains unchanged}. The type system's primary job is to
conservatively track poison propagation through qualified types.

\subsubsection{Type Relationships}

The three type levels are related through systematic transformations:

\begin{itemize}
    \item \textbf{Surface $\to$ Runtime:} Subtyping relation \texttt{ST <: T} (surface type erasure) adds region and closure information
    \item \textbf{Runtime $\to$ Qualified:} Wrapping with poison status and predicates: \texttt{T $\triangleright$ p \{P\}}
    \item \textbf{Function definitions:} Surface types in user syntax become enriched types in type checker, then qualified types in poison-propagating environments
    \item \textbf{Function applications:} Function types store surface types for interface, use typing closures for implementation
\end{itemize}

\textbf{Design Rationale:} This three-level hierarchy cleanly separates the user
interface (surface types), the implementation mechanism (enriched types), and the
analysis results (qualified types). Users write simple, familiar types while the
system performs sophisticated poison tracking and higher-order reasoning behind
the scenes.

\subsection*{Key features}
\begin{itemize}
    \item \textbf{References}: \texttt{ref e} creates, \texttt{e1 := e2} assigns, \texttt{!e} dereferences (memory stores only integers)
    \item \textbf{Functions}: \texttt{fn x : ST1 => e : ST2} creates functions with explicit parameter and return type annotations, \texttt{e1 e2} applies, with closures
    \item \textbf{Pure functions}: \texttt{pfn x : ST1 => pe : ST2} creates pure functions with explicit parameter and return type annotations that cannot have side effects (note: no nested pure functions allowed)
    \item \textbf{Let bindings}: \texttt{let x = e1 in e2} uses type inference for bindings
    \item \textbf{Memory access only through references}: You need a valid
    reference to access memory location. There is no pointer arithmetic (but
    there is aliasing).
    \item \textbf{Qualified types}: All values have qualified types \texttt{Q}
    that combine enriched types with poison status and optional predicate.
    \item \textbf{Predicate-based de-poisoning}: Guards allow selective de-poisoning when runtime conditions are met
    \item \textbf{Sequencing}: \texttt{e1 ; e2} evaluates \texttt{e1} first (which can have any type), discards its value, then evaluates \texttt{e2}
\end{itemize}

Enriched type annotations like region annotations (\texttt{@l}) for references and typing closures
(\texttt{typing\_closure tcl}) for functions are essential to the formal system,
though not part of the user-visible language syntax. The \texttt{l} denotes a memory
location, and \texttt{tcl} represents the typing closure containing the function
parameter name, body, and captured variable typing context. The poison
\texttt{p} represents the function's own poison status (from pc at definition
time). See the formal specification for complete details. Similarly,
qualified types are essential to the formal system, though not part of the
language syntax.

\subsection{Pure Expressions}

The language includes a syntactic subset called \emph{pure expressions} (\texttt{pe}) that exclude side-effecting operations. Pure expressions cannot:
\begin{itemize}
    \item Create new references (\texttt{ref})
    \item Perform assignments (\texttt{:=})
    \item Use \texttt{while} loops. (strictly speaking while loops could be
    side-effects free, but we chose not include them in the ``pure'' subset for
    simplicity. This design choice is motivated by intended use of pure
    functions as type predicates which are likely to be simple expressions
    without complicated control flow.)
    \item Define functions, including pure functions (\texttt{fn} or \texttt{pfn})
    \item Use sequencing (\texttt{;})
    \item Use non-deterministic primitives (\texttt{rand}, \texttt{frand})
    \item Perform I/O (\texttt{printInt}, \texttt{printString})
\end{itemize}

Pure expressions can still dereference existing references, perform arithmetic,
use conditionals, and apply existing pure functions. 

\textbf{Design rationale}: Pure expressions are primarily used within predicates
for de-poisoning. By prohibiting function definition within pure expressions, we
ensure predicates remain simple test expressions rather than complex
computations. The absence of sequencing is natural since pure expressions have
no side effects, making \texttt{pe1 ; pe2} equivalent to just \texttt{pe2}.
However, pure expressions can still apply pre-existing pure functions, enabling
code reuse and composition of predicates.

The \texttt{pfn} construct creates pure functions whose bodies must be pure
expressions, ensuring they cannot introduce side effects. These pure functions
must be defined at the expression level, not within predicates themselves.

\subsubsection{No Nested Pure Functions}

An important consequence of our design is that pure functions cannot define nested
pure functions. Since \texttt{pfn x : ST1 => pe : ST2} requires a pure expression body, and
pure expressions cannot contain \texttt{pfn}, the following patterns are not
possible:

\begin{lstlisting}[language=poison]
(* Not allowed: nested pure function definition *)
pfn outer : int => pfn inner : int => inner + 1 : int : (int -> int)

(* Not allowed: higher-order pure function returning another function *)
pfn make_adder : int => pfn n : int => pfn x : int => x + n : int : (int -> int) : (int -> (int -> int))
\end{lstlisting}

This restriction keeps pure functions \emph{flat} and simple. While it prevents
some functional programming patterns, it aligns with our goal of maintaining
simple, predictable predicates. Pure functions can still:
\begin{itemize}
    \item Call other pre-defined pure functions
    \item Be passed as arguments (through variables)
    \item Be composed through application
\end{itemize}

For example, function composition is still possible through explicit application:
\begin{lstlisting}[language=poison]
pfn add_one : int => x + 1 : int
pfn add_two : int => add_one (add_one x) : int  (* Composition by application *)
\end{lstlisting}

This design choice prioritises simplicity and clarity over full functional
expressiveness, which is appropriate given that pure functions primarily serve
as building blocks for predicates in our fault-tracking type system.

\section{Poison Tracking}

\subsection{Basic Concept}

Values can be \textbf{clean} (reliable) or \textbf{poisoned} (unreliable due to
faults). Any dataflow from a poisoned value is poisoned. 

\begin{lstlisting}[language=poison]
(* Assume get_sensor is a built-in or external function *)
let reading : int = get_sensor unit in  (* May be poisoned due to sensor fault *)
let adjusted : int = reading + 10 in    (* Becomes poisoned if reading is poisoned *)
adjusted
\end{lstlisting}

However, the language also supports \textbf{predicate-based de-poisoning}, where
values can be conditionally cleaned based on runtime tests (discussed in more
detail in Section~\ref{sec:de-poisoning}).

\subsection{Poison Lattice}

\begin{align*}
\text{clean} \sqcup \text{clean} &= \text{clean} \\
\text{clean} \sqcup \text{poisoned} &= \text{poisoned} \\
\text{poisoned} \sqcup \text{poisoned} &= \text{poisoned}
\end{align*}

Conservative: may mark reliable data as poisoned, never marks unreliable data as clean.

\subsection{Predicate-based de-poisoning}
\label{sec:de-poisoning}

The language includes a predicate system that allows selective de-poisoning based on runtime conditions. Predicates come in three forms and are associated with guard tags:

\begin{itemize}
    \item \textbf{Qualified types}: Values have types $T\,\triangleright\,p\,\{P\}$
    where \texttt{T} is the enriched type, \texttt{p} is the poison status, and
    \texttt{P} is an optional predicate.
    \item \textbf{Guard predicates with closure}: \texttt{guard k pred\_closure (x, pe, r)}
    creates a predicate with tag \texttt{k} and pure expression test \texttt{pe}
    (pure function from \texttt{T} to \texttt{bool}). The \texttt{r} captures
    value of variables used in \texttt{pe} at definition time. These enable
    syntactic matching for conditional de-poisoning.
    \item \textbf{Guard-only predicates}: \texttt{guard k} associates a guard tag
    \texttt{k} with a value without a syntactically matchable predicate closure.
    These can be de-poisoned when the tag is known to be true in the logical
    environment, but cannot be matched syntactically via \texttt{match\_guard}.
    \item \textbf{Conditional de-poisoning}: When an \texttt{if} condition
    syntactically matches a guard's test (via \texttt{match\_guard}), values are
    de-poisoned in the true branch and poisoned in the false branch
    \item \textbf{Logical environment}: Environment $\mathcal{L}$ tracks guard
    valuations (\texttt{k : true/false}) to enable predicate evaluation for both
    guard closure and guard-only predicates
\end{itemize}

\textbf{Design Note:} The predicate closure \texttt{pred\_closure (x, pe, r)}
does not explicitly store the expected argument type for the predicate. This
design is sound because predicate matching via \texttt{match\_guard} only occurs
against conditions that have already passed type checking. Since syntactic
matching requires exact structural correspondence with variable substitution,
any predicate that successfully matches a well-typed condition is guaranteed to
be type-compatible with that condition. This approach avoids redundant type
information in predicates while maintaining type safety through the existing
type system.

\textbf{Tag Uniqueness:} The \texttt{match\_guard} function expects to find at
most one matching predicate for any given condition. Predicates parametrised by
their argument will naturally match at most once due to syntactic substitution
differences. However, for predicates that ignore their argument, extra care must
be taken to ensure no two distinct tags are assigned to identical predicate
bodies, as this would cause multiple-match errors in \texttt{match\_guard}.

\textbf{Guard-only limitations:} Guard-only predicates (\texttt{guard k} without
a closure) cannot be matched via \texttt{match\_guard} because they lack the
predicate body required for syntactic comparison. These predicates can only be
de-poisoned when the guard tag \texttt{k} is already known to be true in the
logical environment $\mathcal{L}$ (e.g., from a previous matched guard in an
outer scope or from a while loop guard). This makes guard-only predicates
suitable for propagating guard information without enabling new syntactic
matches.

\textbf{Poisoned references and de-poisoning} Any value obtained from a poisoned
reference is poisoned and could not be de-poisoned further via predicate-based
de-poisoning. To implement this, there is a special typing rule
\texttt{deref\_poisoned} for dereferencing a poisoned reference, which not only
returns a poisoned result but also strips the predicate from its type so it can
no longer be de-poisoned later.

\textbf{Example:}

\begin{lstlisting}[language=poison]
let sensor_val : int = get_sensor unit in   (* Type: int :> poisoned {guard "range" ...} *)
if sensor_val >= 0 then
    sensor_val + 10                  (* De-poisoned: int :> clean *)
else  
    0                                (* Alternative when predicate fails *)
\end{lstlisting}

In this example, if the runtime value of \texttt{sensor\_val} satisfies the range check (\texttt{>= 0}), it becomes de-poisoned in the \texttt{then} branch, allowing safe computation.

\subsection{Examples}

\subsubsection*{Explicit flow} (data dependency)
\begin{lstlisting}[language=poison]
let faulty : int = get_sensor unit in     (* poisoned *)
let result : int = faulty + 1 in          (* poisoned (explicit flow) *)
result
\end{lstlisting}

\subsubsection*{Implicit flow} (control dependency)
\begin{lstlisting}[language=poison]
let faulty : int = get_sensor unit in     (* poisoned *)
if faulty > 0 then 
    42  (* poisoned (implicit flow via pc) *)
else 
    17  (* poisoned (implicit flow via pc) *)
\end{lstlisting}

\subsubsection*{Reference operations}
\begin{lstlisting}[language=poison]
let r : ref int = ref 0 in        (* Clean reference, clean data *)
r := get_sensor unit;             (* Data becomes poisoned, reference stays clean *)
!r                                (* Result: poisoned (from data) *)
\end{lstlisting}

\subsubsection*{Poisoned references}
\begin{lstlisting}[language=poison]
let bad_ref : ref int = get_corrupted_ptr unit in  (* Poisoned reference *)
!bad_ref;      (* Result: poisoned (untrusted reference) *)
bad_ref := 42  (* Poisons data at unknown location *)
\end{lstlisting}

\subsubsection*{Function calls and poisoned function selection}
\begin{lstlisting}[language=poison]
let poisoned_cond : int = get_sensor unit in     (* poisoned *)
let f : int -> int = if poisoned_cond > 0 then   (* pc becomes poisoned *)
    fn x : int => x + 1 : int             (* Function itself gets poisoned *)
else
    fn x : int => x * 2 : int             (* This function also gets poisoned *)
in
let result : int = f 10 in                        (* Application result is poisoned *)
result
\end{lstlisting}

In this example, the function \texttt{f} definition is based on poisoned
condition. The function type captures its own poison status (from pc at
definition time), so both branches produce functions with poisoned status. When
applied, the result inherits this poison. 

\section{Poison Propagation Rules}

\subsection{Core Principles}

\begin{enumerate}
    \item \textbf{Program counter poison} (\texttt{pc}): Tracks when control flow depends on poisoned data
    \item \textbf{Conservative approximation}: All potentially affected computations become poisoned
    \item \textbf{Reference distinction}: Separate tracking for reference reliability vs. data reliability
    \item \textbf{Predicate-aware propagation}: When predicates are available, conditional de-poisoning can occur in if-then-else constructs
    \item \textbf{Guard evaluation}: Environment $\mathcal{L}$ tracks predicate guard evaluations to enable precise poison analysis
\end{enumerate}

\subsection{Propagation Summary}

\begin{itemize}
    \item \textbf{Arithmetic/Logic}: $e_1 \text{ op } e_2 \rightarrow \text{pc} \sqcup \text{poison}(e_1) \sqcup \text{poison}(e_2)$
    \item \textbf{Conditionals}: Both branches get $\text{pc} \sqcup \text{poison}(\text{condition})$, but predicate-based rules can enable de-poisoning
    \item \textbf{Predicate conditionals}: When $\text{match\_guard}(e_1, \Gamma, \sigma) = k$, branches are typed with $\mathcal{L}[k \mapsto \text{true}]$ and $\mathcal{L}[k \mapsto \text{false}]$
    \item \textbf{Variable access}: De-poisoning when guard evaluates to true, definite poisoning when false, standard propagation when guard unknown
    \item \textbf{Assignment}: Operation reliability = $\text{pc} \sqcup \text{poison}(\text{ref})$, Data = $\text{pc} \sqcup \text{poison}(\text{ref}) \sqcup \text{poison}(\text{value})$, invalidates $\mathcal{L}$ entries dependent on written location
    \item \textbf{Dereference}: $\text{pc} \sqcup \text{poison}(\text{ref}) \sqcup \text{poison}(\text{data\_at\_location})$, with predicate-based de-poisoning support
    \item \textbf{Let bindings}: Variable gets $\text{pc} \sqcup \text{poison}(\text{bound\_value})$ and preserves predicates
    \item \textbf{Functions}: Abstraction gets $\text{pc}$ and captures body and
    context, typechecks body to derive return type, application re-typechecks body
    with actual argument predicates and current $\mathcal{L}$
\end{itemize}

\section{Operational Semantics}

\begin{itemize}
    \item $\pi$: Pure-mode flag (\textit{pure} or \textit{impure})
    \item $\mu$: Memory mapping locations to integers (only integers can be stored in memory)
    \item $\rho$: Variable environment mapping variables to values
\end{itemize}

The semantics :
\begin{itemize}
  \item \textbf{Big-step} (natural) semantics that evaluate expressions directly to values and a final configuration; we write informally $\langle \pi, e, \mu, \rho \rangle \Downarrow \langle \mu_1, v \rangle$.
\end{itemize}


\subsection*{Key rules:}


Note that closures $\text{closure}(x, e, \rho)$ and pure closures $\text{pe\_closure}(x,
pe, \rho)$ are runtime values that only appear during execution; they are not
part of the source language and do not have typing rules.



\section{Type System (Typestate)}

\subsection{Environments}

\begin{itemize}
    \item $\Gamma$: Memory typing environment $l : Q$ (memory location qualified types)
    \item $\sigma$: Variable typing environment $x : Q$ (variable qualified types)
    \item $\mathcal{L}$: Logical environment $k : bool \times \text{deps}$ (predicate guard evaluations with dependency tracking)
    \item $\text{pc}$: Program counter poison (current context reliability)
\end{itemize}

Note that both $\Gamma$ and $\sigma$ store \emph{qualified types} $Q$ rather
than simple types, and can include predicates that enable conditional
de-poisoning.

\subsection{Dependency Tracking and Invalidation}

To ensure type system soundness, the logical environment $\mathcal{L}$ tracks
dependencies alongside guard valuations. This prevents unsound de-poisoning when
predicate dependencies are modified:

\begin{itemize}
    \item \textbf{Dependency extraction}: When a guard $k$ is introduced via \texttt{match\_guard}, 
    the system computes a conservative set of heap locations that the predicate may read from
    \item \textbf{Dependency storage}: Guard valuations are stored as $k : bool \times \text{deps}$ 
    where \texttt{deps} is a set ofthe heap locations the predicate depends on
    \item \textbf{Invalidation on write}: When writing to location $l$, all entries in $\mathcal{L}$ 
    whose dependency sets contain $l$ are removed from $\mathcal{L}$
    \item \textbf{Conservative analysis}: Dependency extraction uses syntactic analysis of 
    predicate expressions to identify heap reads through variables and dereferences
\end{itemize}

This mechanism ensures that predicate-based de-poisoning remains sound even when
the heap locations that predicates depend on are modified between guard introduction
and consumption.

\subsection{Judgment Form}

\[
\pi \; ;\; \text{pc} \; ;\; \Gamma \; ;\; \sigma \; ;\; \mathcal{L} \vdash e : Q \dashv \Gamma_1 ; \mathcal{L}_1
\]

Reads: "In pure mode $\pi$ and poison context $\text{pc}$, with memory typing environment $\Gamma$,
variable typing environment $\sigma$, and logical environment $\mathcal{L}$,
expression $e$ has qualified type $Q$, producing updated memory typing
environment $\Gamma_1$ and updated logical environment $\mathcal{L}_1$."

The qualified type $Q$ encapsulates both the base type, poison status, and any
associated predicates. The logical environment $\mathcal{L}$ tracks the current
valuation of predicate guards, enabling predicate-based typing decisions. The
output logical environment $\mathcal{L}_1$ reflects any updates to guard
valuations learned during typing (e.g., after a while loop terminates) or
because of invalidation. due to changed dependencies.

\subsection{Key Typing Rules}

\subsubsection*{References with region types:}
\[
\frac{\Gamma(l) = Q}{\pi \; ;\; \text{pc} \; ;\; \Gamma \; ;\; \sigma \; ;\; \mathcal{L} \vdash l : \text{ref } T @ l \triangleright \text{clean} \{ \text{no\_pred} \} \dashv \Gamma}
\]

This rule type actual memory location value, which is always clean.

\subsubsection*{Variable access with predicate support:}
When variables have associated predicates, the typing rule checks the logical environment $\mathcal{L}$ to determine if de-poisoning can occur. Both guard closures and guard-only predicates behave identically when checking $\mathcal{L}$:

\[
\frac{\sigma(x) = T \triangleright p \{ \text{guard } k \text{ [pred\_closure] } (y, pe, r) \} \quad \mathcal{L}(k) = \text{true}}{\pi \; ;\; \text{pc} \; ;\; \Gamma \; ;\; \sigma \; ;\; \mathcal{L} \vdash x : T \triangleright \text{pc} \{ \text{no\_pred} \} \dashv \Gamma}
\]

\[
\frac{\sigma(x) = T \triangleright p \{ \text{guard } k \text{ [pred\_closure] } (y, pe, r) \} \quad \mathcal{L}(k) = \text{false}}{\pi \; ;\; \text{pc} \; ;\; \Gamma \; ;\; \sigma \; ;\; \mathcal{L} \vdash x : T \triangleright \text{poisoned} \{ \text{no\_pred} \} \dashv \Gamma}
\]

where \texttt{[pred\_closure]} denotes that the predicate may be either \texttt{guard k pred\_closure (y, pe, r)} or simply \texttt{guard k}. The presence or absence of the closure does not affect the typing behaviour when the guard tag is known in $\mathcal{L}$.

\subsubsection*{Reference creation:}
\[
\frac{\pi \; ;\; \text{pc} \; ;\; \Gamma \; ;\; \sigma \vdash e : T \rhd p \dashv \Gamma_1 \quad \text{fresh } l}
{\pi \; ;\; \text{pc} \; ;\; \Gamma \; ;\; \sigma \vdash \text{ref } e : \text{ref } T @ l \rhd \text{pc} \dashv \Gamma_1[l \mapsto T \triangleright \text{pc} \sqcup p]}
\]

\subsubsection*{Assignment:}
\[
\frac{\pi \; ;\; \text{pc} \; ;\; \Gamma \; ;\; \sigma \vdash e_1 : \text{ref } T @ l \rhd p_1 \dashv \Gamma_1 \quad \pi \; ;\; \text{pc} \; ;\; \Gamma_1 \; ;\; \sigma \vdash e_2 : T \rhd p_2 \dashv \Gamma_2}
{\pi \; ;\; \text{pc} \; ;\; \Gamma \; ;\; \sigma \vdash e_1 := e_2 : \text{unit} \rhd \text{pc} \sqcup p_1 \dashv \Gamma_2[l \mapsto T \triangleright \text{pc} \sqcup p_1 \sqcup p_2]}
\]

\subsubsection*{Dereference:}
\[
\frac{\pi \; ;\; \text{pc} \; ;\; \Gamma \; ;\; \sigma \vdash e : \text{ref } T @ l \rhd p_1 \dashv \Gamma_1 \quad \Gamma_1(l) = T \triangleright p}
{\pi \; ;\; \text{pc} \; ;\; \Gamma \; ;\; \sigma \vdash !e : T \rhd \text{pc} \sqcup p_1 \sqcup p \dashv \Gamma_1}
\]

\subsubsection*{Function abstraction:}
Function types include a \emph{typing closure} that captures the function
body and variable typing context at definition time. Both the parameter
type $ST_1$ and return type $ST_2$ are explicitly provided as surface type annotations.
The function body is typechecked at definition time to verify it produces a type compatible with the declared return type $ST_2$.
\[
\frac{\begin{array}{c}
\pi \; ;\; \text{pc} \; ;\; \Gamma \; ;\; \sigma[x \mapsto T_1 \triangleright \text{pc} \{ \text{no\_pred} \}] \; ;\; \mathcal{L} \vdash e : T_2 \triangleright p_1 \{ P_1 \} \dashv \Gamma_1 ; \mathcal{L}_1 \\
ST_1 <: T_1 \\
ST_2 <: T_2
\end{array}}
{\pi \; ;\; \text{pc} \; ;\; \Gamma \; ;\; \sigma \; ;\; \mathcal{L} \vdash \text{fn } x : ST_1 \Rightarrow e : ST_1 \rightarrow ST_2 \# \text{pc} \; \text{typing\_closure}(x, e, \sigma) \triangleright \text{pc} \{ \text{no\_pred} \} \dashv \Gamma ; \mathcal{L}}
\]

\noindent For pure functions, the enriched type uses a separate pure arrow and
closure: $ST_1 \rightarrow ST_2 \# p\; \text{typing\_pclosure}(x,pe,\sigma)$,
and application uses the pure closure in place of the impure one.

\paragraph{Pure-mode typing.} We thread a pure-mode parameter \(\pi\) through
all typing judgements. Impure function application (\textsc{app\_fn}) is only
available in impure mode, expressed by a conclusion headed with \textit{impure}.
Pure function application (\textsc{app\_pfn}) types the callee’s body in pure
mode but preserves the caller’s \(\pi\) in the conclusion. Thus, pure code can
only call pure functions; impure code can call either. The mode \(\pi\) is
orthogonal to the poison \(pc\): standard poison joins and propagation are
unchanged. We enter programs with \(\pi = \textit{impure}\).

\subsubsection*{Function application:}
The function application typing rule extracts the typing closure from the function type and types the body at application time. This enables predicate-based de-poisoning by using the actual argument's predicates and the current logical environment $\mathcal{L}$. The subtyping checks $ST_1 \!<:\! T_1$ and $ST_2 \!<:\! T_2$ ensure the surface parameter and return types are compatible with the actual argument and body types respectively.

\[
\frac{
\begin{array}{c}
\text{pc} ; \Gamma ; \sigma ; \mathcal{L} \vdash e_1 : ST_1 \rightarrow ST_2 \# p \;\; \text{typing\_closure}(x, e, \sigma_1) \triangleright p_1 \{ P_1 \} \dashv \Gamma_1 ; \mathcal{L}_1 \\
\text{pc} ; \Gamma_1 ; \sigma ; \mathcal{L}_1 \vdash e_2 : T_1 \triangleright p_2 \{ P_2 \} \dashv \Gamma_2 ; \mathcal{L}_2 \\
ST_1 <: T_1 \\
\text{join } p_1 \; pc ; \Gamma_2 ; \sigma_1[x \mapsto T_1 \triangleright \text{pc} \sqcup p_2 \{ P_2 \}] ; \mathcal{L}_2 \vdash e : T_2 \triangleright p_3 \{ P_3 \} \dashv \Gamma_3 ; \mathcal{L}_3 \\
ST_2 <: T_2
\end{array}
}{
\text{pc} ; \Gamma ; \sigma ; \mathcal{L} \vdash e_1 \; e_2 : T_2 \triangleright \text{pc} \sqcup p_3 \{ P_3 \} \dashv \Gamma_3 ; \mathcal{L}_3
}
\]


Key aspects of function application:
\begin{itemize}
    \item The typing closure $(x, e, \sigma_1)$ contains the parameter name, function body, and captured variable context
    \item The function body $e$ is typed using the captured context $\sigma_1$ extended with the actual argument's type
    \item The argument's qualified type includes its predicates $P_2$, enabling de-poisoning in the function body
    \item The current logical environment $\mathcal{L}_2$ is used, allowing guard-based de-poisoning
    \item The result poison $p_3$ reflects any de-poisoning that occurred in the function body
    \item Function types store both surface parameter type $ST_1$ and return type $ST_2$, with body typing occurring at both definition and application time
\end{itemize}


Example of predicate-based de-poisoning in functions:
\begin{lstlisting}[language=poison]
(* Assume get_sensor is a built-in that returns poisoned sensor readings with predicates *)
let sensor1 : int = get_sensor unit in  (* Type: int :> poisoned {guard "temp" (y, y >= 0, _)} *)
let sensor2 : int = get_sensor unit in  (* Type: int :> poisoned {guard "pressure" (z, z >= 0, _)} *)

let f : int -> int = fn x : int =>      (* Captures sensor1 with its predicate *)
    if sensor1 >= 0 then                (* Matches guard "temp" test expression *)
        if x >= 0 then                  (* Matches guard "pressure" test expression *)
            sensor1 + x                 (* Both de-poisoned: clean result *)
        else
            -1                          (* x still poisoned *)
    else
        -2                              (* sensor1 still poisoned *)
    : int                               (* Explicit return type annotation *)
in
let result : int = f sensor2 in         (* Application types body with predicates *)
result                                  (* Result: int :> clean *)
\end{lstlisting}

In this example:
\begin{itemize}
    \item \texttt{sensor1} has predicate \texttt{guard "temp" pred\_closure (y, y >= 0, \_)} - a test that checks if the value is non-negative
    \item \texttt{sensor2} has predicate \texttt{guard "pressure" pred\_closure (z, z >= 0, \_)} - also tests for non-negative
    \item When \texttt{f} is applied to \texttt{sensor2}, the function body is typed with:
    \begin{itemize}
        \item Captured context $\sigma_1$ containing \texttt{sensor1} with its predicate
        \item Parameter \texttt{x} bound to \texttt{sensor2}'s type with its predicate  
        \item Current logical environment $\mathcal{L}$ (initially may be empty)
    \end{itemize}
    \item During typing of the function body:
    \begin{itemize}
        \item \texttt{if sensor1 >= 0} syntactically matches the test expression in \texttt{guard "temp"}, so the then branch is typed with $\mathcal{L}[$\texttt{"temp"} $\mapsto \text{true}]$
        \item \texttt{if x >= 0} syntactically matches the test expression in \texttt{guard "pressure"}, so its then branch is typed with $\mathcal{L}[$\texttt{"temp"} $\mapsto \text{true}, $\texttt{"pressure"} $\mapsto \text{true}]$
        \item In the innermost branch, both values are de-poisoned because their guards are true in $\mathcal{L}$
        \item The result \texttt{sensor1 + x} has type \texttt{int} $\triangleright$ \texttt{clean}
    \end{itemize}
\end{itemize}

\subsubsection*{Predicate-based conditional typing:}
The type system includes two conditional typing rules that handle predicate-based de-poisoning:

\textbf{Standard conditional rule (when no matching predicate found):}
\[
\frac{\begin{array}{c}
\text{pc} ; \Gamma ; \sigma ; \mathcal{L} \vdash e_1 : \text{bool} \triangleright p_1 \{ P_1 \} \dashv \Gamma_1 \\
\text{join pc } p_1 ; \Gamma_1 ; \sigma ; \mathcal{L} \vdash e_2 : T \triangleright p_2 \{ P_2 \} \dashv \Gamma_2 \\
\text{join pc } p_1 ; \Gamma_1 ; \sigma ; \mathcal{L} \vdash e_3 : T \triangleright p_3 \{ P_3 \} \dashv \Gamma_3
\end{array}}
{\text{pc} ; \Gamma ; \sigma ; \mathcal{L} \vdash \text{if } e_1 \text{ then } e_2 \text{ else } e_3 : T \triangleright \text{join pc } p_1 p_2 p_3 \{ \text{no\_pred} \} \dashv \text{JoinG } \Gamma_2 \Gamma_3}
\]

\textbf{Predicate-based conditional rule (when matching predicate found):}
\[
\frac{\begin{array}{c}
\text{pc} ; \Gamma ; \sigma ; \mathcal{L} \vdash e_1 : \text{bool} \triangleright p_1 \{ P_1 \} \dashv \Gamma_1 ; \mathcal{L}_1 \\
\text{match\_guard}(e_1, \Gamma, \sigma) = k \\
pc ; \Gamma_1 ; \sigma ; \mathcal{L}_1[k \mapsto \text{true}] \vdash e_2 : T \triangleright p_2 \{ P_2 \} \dashv \Gamma_2 ; \mathcal{L}_2 \\
pc ; \Gamma_1 ; \sigma ; \mathcal{L}_1[k \mapsto \text{false}] \vdash e_3 : T \triangleright p_3 \{ P_3 \} \dashv \Gamma_3 ; \mathcal{L}_3
\end{array}}
{\text{pc} ; \Gamma ; \sigma ; \mathcal{L} \vdash \text{if } e_1 \text{ then } e_2 \text{ else } e_3 : T \triangleright \text{join pc } p_2 p_3 \{ \text{no\_pred} \} \dashv \text{JoinG } \Gamma_2 \Gamma_3 ; \mathcal{L}_1}
\]

The key difference is that when $\text{match\_guard}(e_1, \Gamma, \sigma) = k$, the branches are typed with different valuations of guard $k$ in environment $\mathcal{L}$. This enables:
\begin{itemize}
    \item \textbf{De-poisoning in then branch}: Values guarded by $k$ become clean when $k$ evaluates to true
    \item \textbf{Definite poisoning in else branch}: Values guarded by $k$ become poisoned when $k$ evaluates to false
    \item \textbf{Precise fault tracking}: The type system can track exactly which runtime conditions affect poison status
\end{itemize}

NOTE: In de-poisoning case when predicate is matched we do not propage poison from condition check to the branches.

\paragraph{Note on $\mathcal{L}_1$ propagation:} The predicate-based conditional rule returns $\mathcal{L}_1$ (the environment after evaluating $e_1$) rather than attempting to join $\mathcal{L}_2$ and $\mathcal{L}_3$ from the branches. This is a conservative but sound design choice, as $\mathcal{L}_2$ contains $k \mapsto \text{true}$ while $\mathcal{L}_3$ contains $k \mapsto \text{false}$, making them contradictory. After the conditional completes, we don't know which branch was taken at runtime, so we revert $k$ to unknown status by returning $\mathcal{L}_1$.

\paragraph{Alternative design consideration:} A more precise approach would implement a proper join operation for logical environments where:
\begin{itemize}
    \item Guards with the same value in both branches are preserved
    \item Guards with contradictory values are removed (become unknown)
    \item Guards only present in one branch are removed
\end{itemize}
This would provide more precision but adds complexity. The current approach is simpler while still enabling the key benefit: predicate-based de-poisoning within each branch of the conditional.

\subsubsection*{While loops}

While loop typing is flow-sensitive, following Hunt \& Sands 2006~\cite{DBLP:conf/popl/HuntS06}. We index by
iteration $i$ the program-counter poison $p_i$, the memory typing environment $\Gamma_i$, and the
logical environment $\mathcal{L}_i$. The least fixpoint is reached at the smallest $n$ such that
\(p_{n+1}=p_n\), \(\Gamma_{n+1}=\Gamma_n\), and \(\mathcal{L}_{n+1}=\mathcal{L}_n\).

The formal specification provides \emph{two} mutually exclusive typing rules for loops:
\begin{itemize}
  \item \textbf{\texttt{while\_depoison}}~-- applied when
        \texttt{match\_guard}$(e_1,\Gamma,\sigma)=k$ detects that the loop
        guard $e_1$ syntactically matches a predicate guard tagged $k$.
        The guard expression $e_1$ is typed as a \emph{pure} expression (no effects),
        so it does not modify $\Gamma$ or $\mathcal{L}$. The body $e_2$ is typed under
        the logical environment $\mathcal{L}[k\mapsto{}\text{true}]$, thereby treating
        every value guarded by $k$ as \emph{clean} within the body. Across iterations we compute
        the recurrences
        \[
          p_{i+1} = p_i \sqcup p_{2,i},\qquad
          \Gamma_{i+1} = \text{JoinG}(\Gamma_i,\Gamma_{2,i}),\qquad
          \mathcal{L}_{i+1} = \text{JoinL}(\mathcal{L}_i,\mathcal{L}_{2,i}),
        \]
        where $(\cdot)_{2,i}$ denotes the quantities produced by the $i$-th typing of the body.
        At loop exit we just observed $e_1$ evaluate to \emph{false}, hence the output logical
        environment sets $k\mapsto{}\text{false}$ (with the appropriate dependency set).

  \item \textbf{\texttt{while\_std}}~-- the conservative fallback when no
        suitable guard is found.  Here the guard $e_1$ may be impure and can modify both
        $\Gamma$ and $\mathcal{L}$. The same fixpoint recurrences are used with the
        \texttt{JoinG} and \texttt{JoinL} operators.
\end{itemize}

\paragraph{Join on the logical environment.} The loop back-edge uses an \emph{agreement-only} join on
$\mathcal{L}$ to prevent iteration-local facts from leaking:
for each tag $k$,
\[
  \text{JoinL}(\mathcal{L}_1,\mathcal{L}_2)(k) =
  \begin{cases}
    (b,\; \text{deps}_1 \cup \text{deps}_2), & \text{if } \mathcal{L}_1(k)=(b,\text{deps}_1) \text{ and } \mathcal{L}_2(k)=(b,\text{deps}_2); \\
    \text{unknown}, & \text{otherwise (absent on one side or booleans disagree).}
  \end{cases}
\]
This definition is idempotent, commutative, and associative, composes with dependency invalidation
on writes, and ensures the sequence $\mathcal{L}_{i+1}=\text{JoinL}(\mathcal{L}_i,\mathcal{L}_{2,i})$ is
descending in information and stabilises. Using a union that preserves single-sided or contradictory
facts would be unsound for de-poisoning.

\paragraph{Precision gained.} By exploiting syntactic matching against guard
predicates, the \texttt{while\_depoison} rule can identify cases where a loop
condition constitutes a run-time validity check.  Inside the loop body the
associated data are thus considered reliable, enabling more precise reasoning
about fault propagation compared with the standard rule.

\paragraph{Loop termination handling.} When a loop terminates we have just observed
the guard expression $e_1$ evaluate to \emph{false}.  Hence the matching tag
$k$ is \textbf{known to be false at the loop exit}, and every value whose
qualified type is guarded by $k$ ought to be marked \emph{poisoned}. The
typing rule correctly updates the logical environment $\mathcal{L}$
by setting $k$ to false in the output environment (i.e., $\mathcal{L}_n[k \mapsto \text{false}]$ with
merged dependencies), 
allowing post-loop code to properly classify values guarded by $k$ as poisoned. 
This provides precise fault tracking both during and after loop execution.

\subsection{Region Types for Soundness}

Region types $\text{ref } T @ l$ solve the soundness problem in dereference:
\begin{itemize}
    \item \textbf{Problem}: Without regions, $!e$ rule couldn't determine which location to check in $\Gamma$ for poison status
    \item \textbf{Solution}: Type $\text{ref } T @ l$ explicitly encodes the target location $l$
    \item \textbf{Benefit}: Statically tracks which specific location each reference points to
\end{itemize}

\subsection{Environment Evolution}

Memory typing environment $\Gamma$ evolves through:
\begin{itemize}
    \item \textbf{Reference creation}: Extends $\Gamma$ with both type and poison status
    \item \textbf{Assignment}: Updates poison status of data in $\Gamma$
    \item \textbf{Control flow}: Joins environments from different branches
    \item \textbf{Loops}: Iterative fixpoint computation for flow-sensitive analysis
\end{itemize}

Logical environment $\mathcal{L}$ evolves through:
\begin{itemize}
    \item \textbf{Sequence}: Chained through expressions - first expression produces $\mathcal{L}_1$, second uses $\mathcal{L}_1$ and produces $\mathcal{L}_2$
    \item \textbf{While loops}: After a predicate-matched loop terminates, the guard tag is set to false in $\mathcal{L}_1$
    \item \textbf{Assignment}: Invalidates entries whose dependencies include the written location
    \item \textbf{Most other constructs}: Pass through unchanged
\end{itemize}

\section{Formal Properties}

\subsection{Non-interference}

We define the non-interference property for the Poison language, which ensures
that the program's observable behaviour depends solely on critical inputs.
Additionally, the property allows for fault isolation, where some critical
inputs must be checked for validity before use.

\paragraph{We distinguish three classes of inputs:}
\begin{itemize}
  \item \textbf{Critical inputs}: exactly those listed in the specification $S$;
  they determine the guard regime and thus govern non-interference.
  \item \textbf{Contextual non-critical inputs}: not listed in $S$ but present
  in the initial typing context because the program refers to them. They are
  always recorded as poisoned with \texttt{no\_pred} in the context and must be
  instantiated at runtime to avoid stuckness; they cannot be de-poisoned.
  \item \textbf{Unused inputs}: neither listed in $S$ nor present in the initial
  typing context; they are irrelevant and may be absent at runtime.
\end{itemize}

\begin{definition}[Input/Output Classification]
Inputs consist of initial memory and variables:
\begin{itemize}
    \item Initial memory: $\mu_0 : \text{Location} \rightarrow \text{Value}$
    \item Initial variables: $\rho_0 : \text{Variable} \rightarrow \text{Value}$
\end{itemize}
We fix a \emph{critical-input specification} $S$ comprising lists of critical entries with their qualified types:
\[
  S.\mathrm{locs} = \bigl[(\ell_1, Q_1),\ldots,(\ell_m, Q_m)\bigr],\qquad
  S.\mathrm{vars} = \bigl[(x_1, Q'_1),\ldots,(x_n, Q'_n)\bigr].
\]
Entries not listed in $S$ are considered non-critical and are \emph{unconstrained by the specification}.

Outputs consist of:
\begin{itemize}
    \item Final memory locations: $\mu_f : \text{Location} \rightarrow \text{Value}$
    \item Program return value: $v$
\end{itemize}
With observable outputs designated as $(L_{obs}, v_{obs})$ where $v_{obs} \in \{\text{true}, \text{false}\}$ indicates whether the return value is observable.
Non-observable outputs are unconstrained.

Note: Variable bindings are local to expressions and do not persist as program outputs. Only memory locations and the return value constitute observable outputs.
\end{definition}

\begin{definition}[Guard Regime]
Let $K_S$ be the finite set of guard tags that occur in predicates attached to
entries of $S$.

A \emph{guard regime} is a total mapping $R : K_S \to
\{\text{true},\text{false}\}$ assigning a truth value to each tag. It prescribes
which guards are considered true in an execution.

The specification $S$ determines the constraints for critical inputs and the
initial typing context; a guard regime $R$ supplies truth values for tags in
$K_S$ that instantiate these constraints at comparison time.

\paragraph{Regime/Specification consistency.}
For a single run with initial memory and variable environments $(\mu,\rho)$, we say that $R$ is \emph{consistent} with $(\mu,\rho)$ when, for every critical entry in $S$:
\begin{itemize}
  \item if the entry has no predicate, it is trivially consistent (statically
  valid by $S$);
  \item if it has a \emph{guard-only} tag $k$, there is no consistency
  requirement and it is trivially consistent.
  \item if it has a predicate \emph{with closure} tagged $k$, the predicate
  closure applied to the entry's runtime value evaluates to $\text{true}$ iff
  $R(k)=\text{true}$ and to $\text{false}$ iff $R(k)=\text{false}$;
\end{itemize}

\textbf{Excluding non-critical guards.} Regimes range only over $K_S$. Tags originating solely from non-critical inputs are outside the regime and, by construction, cannot influence the observable behaviour required by the property.
\end{definition}

\begin{definition}[Non-Interference Property]
A program $e$ satisfies non-interference with respect to the critical-input specification $S$ and observable outputs if, for each guard regime $R$:

Given two executions with the same observable-behaviour form as defined above,
$\langle e, \mu_1, \rho_1 \rangle \Downarrow \langle \mu_{1,1}, v_1 \rangle$ and
$\langle e, \mu_2, \rho_2 \rangle \Downarrow \langle \mu_{2,1}, v_2 \rangle$, the following hold:
\begin{itemize}
    \item \textbf{Agreement on critical inputs under $R$:}
    \begin{itemize}
        \item For critical entries without predicates, the corresponding initial
        values are identical across the two executions.
        \item For entries with predicate \emph{closures} tagged $k$: when
        $R(k)=\text{true}$ the initial values satisfy the closure and are
        identical; when $R(k)=\text{false}$ the initial values may differ but
        both fail the closure.
        \item For entries with \emph{guard-only} tag $k$: when $R(k)=\text{true}$ the
        initial values are identical; when $R(k)=\text{false}$ the initial values may
        differ.
    \end{itemize}
    \item \textbf{Non-critical inputs:} unconstrained.
\end{itemize}
Then the observable outputs coincide:
\begin{itemize}
    \item $\forall l \in L_{obs}.\; \mu_{1,1}(l) = \mu_{2,1}(l)$
    \item If $v_{obs} = \text{true}$ then $v_1 = v_2$.
\end{itemize}
\end{definition}

The key insight initial values asymmetry governed by the regime $R$:
\begin{itemize}
    \item If a critical entry is annotated with tag $k$ and $R(k)=\text{true}$,
    it is treated as valid; the compared runs must agree on its initial value
    (for predicate closures), and the observable behaviour may depend on that
    value.
    \item If $R(k)=\text{false}$, the entry is treated as invalid; the compared
    runs need not agree on its initial value (for predicate closures, both must
    fail), yet observables must coincide.
    \item Entries without predicates are statically valid by $S$ and must agree
    across runs.
\end{itemize}

This is a runtime property which is impossible, in general, to guarantee statically. However,
in the next definition, we will show that we can conservatively approximate it
via the type system.

\begin{theorem}[Type System Enforces Non-Interference]
Assume a well-formed critical-only input specification $S$ and observable spec
$\text{obs}$. For any initial typing context $(pc,\Gamma,\sigma,\mathcal{L})$
that matches $S$ (namely: $pc=\text{clean}$, $\mathcal{L}=\emptyset$; the
entries of $S$ appear in $\Gamma$ and $\sigma$ with exactly the specified
qualified types; any other entries present are poisoned with $\text{no\_pred}$),
the following implication holds:

If typechecking $e$ under this context succeeds and the inferred
result satisfies the observable constraints (observable locations and, if
requested, the return are clean), then for \emph{any} abstract predicate 
satisfaction relation $\text{PredSat} : \text{Tag} \to \text{Memory} \to 
\text{Value} \to \text{Prop}$ and dependency relation 
$\text{PredicateDeps} : \text{Tag} \to \text{Deps} \to \text{Prop}$ 
satisfying the axioms (syntactic matching and memory monotonicity as 
described below), the program $e$ satisfies the non-interference
property with respect to $S$, $\text{obs}$, $\text{PredSat}$, and $\text{PredicateDeps}$.

The runtime premises for each compared run are: (i) critical inputs conform to
$S$, and (ii) the concrete initial memory and environment instantiate the whole
context $(pc,\Gamma,\sigma,\mathcal{L})$.
\end{theorem}

The predicate system enables graceful degradation: programs can produce correct
outputs even when some critical inputs are faulty, provided those inputs have
associated validity checks that are properly exercised.

\subsection{Abstract Predicate Satisfaction}

The non-interference proof employs an \emph{abstract predicate satisfaction} approach to handle predicates with captured lexical environments. This design addresses the fundamental architectural constraint that runtime values carry no type information or tags.

\subsubsection{The Challenge}

Predicate closures \texttt{guard k pred\_closure (x, pe, r)} capture a typing environment \texttt{r : List (Name × QTyp)} at definition time. However, the non-interference proof must reason about whether predicates are satisfied by runtime values without access to the captured runtime variable bindings. The typing environment contains only type information, not the actual values of captured variables.

\subsubsection{The Solution: Abstract Satisfaction Relation}

Instead of concretely evaluating predicates, the proof treats predicate satisfaction as an abstract mathematical relation:
\[
\text{PredSat} : \text{Tag} \to \text{Memory} \to \text{Value} \to \text{Prop}
\]

This relation specifies whether a predicate (identified by its tag) is satisfied by a value in a given memory state, without requiring concrete evaluation. Key design aspects:

\begin{itemize}
    \item \textbf{No variable environment parameter}: Captured variables are encoded in the tag's identity. Each tag uniquely identifies both the predicate body and its captured lexical context.
    \item \textbf{Memory parameter}: Required because predicates may dereference heap locations (e.g., \texttt{!x}), which can change during execution.
    \item \textbf{Universal quantification}: The non-interference theorem quantifies over all \texttt{PredSat} relations satisfying certain axioms, proving the property holds regardless of what predicates mean.
\end{itemize}

\subsubsection{Constraining Axioms}

The abstract \texttt{PredSat} relation is constrained by axioms that capture the type system's mechanisms without specifying predicate meanings:

\paragraph{Syntactic Matching Axiom.} If a condition syntactically matches a tagged predicate via \texttt{match\_guard} and evaluates to true, then the predicate is satisfied:
\[
\forall \text{ condition, } k, \text{deps}, \Gamma, \sigma, \mu, \rho, v.\;
\text{matchGuard}(\text{condition}, \Gamma, \sigma) = (k, \text{deps})
\land \text{terminatesTo}(\text{condition}, \mu, \rho, \text{true})
\implies \text{PredSat}\; k\; \mu\; v
\]

This axiom connects the type system's syntactic predicate matching (which returns both a tag and its dependencies) with the abstract satisfaction relation.

\paragraph{Memory Monotonicity Axiom.} Predicates are stable under memory changes that do not affect their dependencies:
\[
\forall k, \mu, \mu', v, \text{deps}.\;
\text{PredicateDeps}\; k\; \text{deps}
\land (\forall l \in \text{deps}.\; \mu(l) = \mu'(l))
\implies (\text{PredSat}\; k\; \mu\; v \iff \text{PredSat}\; k\; \mu'\; v)
\]

This axiom ensures the abstract relation behaves consistently with the type system's dependency tracking and guard invalidation on assignment.

\subsubsection{Relationship to Regimes}

The guard regime $R$ and abstract satisfaction \texttt{PredSat} serve complementary but independent roles:

\begin{itemize}
    \item \textbf{PredSat} (abstract satisfaction): Universally quantified parameter specifying what predicates mean abstractly. Fixed throughout the proof for all regime choices.
    \item \textbf{Regime $R$} (execution path): Existentially chosen mapping recording which predicates were satisfied for specific runtime values in a particular execution. Different regimes represent different program runs.
\end{itemize}

The regime \emph{witnesses} facts about concrete runtime values under the abstract \texttt{PredSat}, rather than defining what \texttt{PredSat} means. Regime consistency requires that for guard closures, $R(k) = \text{true}$ iff \texttt{PredSat} holds for the runtime value, ensuring the regime correctly records predicate satisfaction.

\subsubsection{Benefits}

This abstract approach correctly handles predicates that capture variables,
including shadowing, without evaluating them in incorrect environments.

The theorem thus states: for \emph{any} interpretation of predicates satisfying
the type system's operational constraints (syntactic matching and dependency
tracking), well-typed programs satisfy non-interference.

\subsection{Non-Completeness}

The Poison language's type system cannot achieve completeness due to its
flow-sensitive nature and conservative fault propagation approximation. Tracking
program counter poison ($pc$) and join operations at control flow merge points
leads to over-approximation. Values in a poisoned $pc$ context are marked as
poisoned, even if they are constants or depend on clean values. This is a
trade-off for soundness, preferring to mark clean values as poisoned rather than
miss poisoned ones. Thus, while clean-marked values are soundly clean,
poisoned-marked values are not necessarily poisoned.



The typestate (memory typing environments $\Gamma$) is monotone only with
respect to base types (\texttt{T}): once a location receives a type, that type
is never changed, though its poison status may evolve through assignments and
control flow. New locations may be added to the environment (e.g., through
\texttt{ref} expressions), but existing locations retain their base types. The
poison components (flag and predicate) within $\Gamma$ is permitted to change in
both directions: ordinary computation and control-flow joins may mark additional
entities as poisoned, whereas \emph{sanitisation} constructs could clear poison
from selected data. Preservation therefore guarantees that the analysis remains
sound while accommodating monotone growth of type information and non-monotone
evolution of poison information.



\bibliographystyle{plain}

\nocite{*}
\bibliography{poison-lang}

\end{document} 